\documentclass[11pt]{article}

\usepackage{amsmath,amssymb,amsthm}
\usepackage[margin=1.1in]{geometry}
\usepackage{bm}
\usepackage{booktabs}
\usepackage{array}
\usepackage[colorlinks=true,linkcolor=blue,citecolor=blue]{hyperref}

\newtheorem{proposition}{Proposition}
\newtheorem{remark}{Remark}
\newcommand{\R}{\mathbb{R}}
\newcommand{\bu}{\mathbf{u}}
\newcommand{\bv}{\mathbf{v}}
\newcommand{\bw}{\mathbf{w}}
\newcommand{\bn}{\mathbf{n}}
\newcommand{\bx}{\mathbf{x}}
\newcommand{\bOmega}{\boldsymbol{\Omega}}
\newcommand{\bomega}{\boldsymbol{\omega}}
\newcommand{\dS}{\,\mathrm{d}S}
\newcommand{\dV}{\,\mathrm{d}V}
\newcommand{\dc}{\,\mathrm{d}c}
\newcommand{\Div}{\operatorname{div}}
\newcommand{\code}[1]{\texttt{#1}}
\newcommand{\nuh}{\hat{\boldsymbol{\nu}}}
\newcommand{\Se}{\mathbf{S}_e}

\title{Dahlen's fluid closure and its generalisation:\\
       hydrostatic fluid regions in a static\\
       elasto-gravitational body of any shape}
\author{AdGIA method notes}
\date{}

\begin{document}
\maketitle
\sloppy

\noindent
This note is a self-contained account of one method: Dahlen's static
treatment of a fluid region, in which the fluid's displacement is
eliminated and the fluid enters the solid's problem through its
density and the interface, and the generalisation of that treatment to
fluid regions of arbitrary shape, to solids with arbitrary pre-stress,
and to rotating bodies. The generalisation adds one one-dimensional
unknown per connected fluid region, defined on the level foliation of
the background potential, and couples it to the three-dimensional
problem as a symmetric border. On a spherical body it reduces to
Dahlen's method at every degree $\ell \ge 1$ and to the compressible
degree-0 problem at $\ell = 0$; off the sphere it is, as far as we
know, the only static method for a non-neutrally stratified fluid
region.

Every statement is tagged \emph{derived} (a calculation written out
here), \emph{argued} (a claim whose proof is sketched but not
completed), or \emph{measured} (a numerical fact from this library).
The broader setting --- why the static problem with the fluid kept as
an elastic material is ill posed when the stratification is not
neutral, what dissipation selects, and how the other formulations of
the library stand --- is \code{doc/static\_fluid\_core.tex}; it is
summarised in \S\ref{sec:why} only as far as this note needs it.

\tableofcontents

%=====================================================================
\section{Setting and notation}
\label{sec:setting}

\paragraph{The body.} A body $B \subset \R^d$, $d = 2$ or $3$,
consists of solid regions and fluid regions separated by interfaces.
One fluid region $F$ with its wall $\Sigma = \partial F$ is discussed;
several regions are independent copies. The unit normal $\bn$ on
$\Sigma$ points out of $F$; where the solid's own outward normal
$\mathbf m = -\bn$ is used, the text says so. $G$ is the gravitational
constant.

\paragraph{The background.} The body is in static equilibrium in the
total background potential $\Phi_0$, gravitational plus, for a rotating
body, centrifugal (\S\ref{sec:rotation}). In the fluid the equilibrium
is hydrostatic,
\begin{equation}
  \nabla p^0 = -\rho\,\nabla\Phi_0 \quad\text{in } F .
  \label{eq:hydro}
\end{equation}
Taking the curl, $\nabla\rho\times\nabla\Phi_0 = \mathbf 0$: the
density and the pressure are constant on each connected component of
each level set of $\Phi_0$ within $F$, and $F$ carries the \emph{level
foliation} of $\Phi_0$. We write $\rho(\Phi_0)$, $\rho_\Phi =
\mathrm d\rho/\mathrm d\Phi_0$ (per component), $g = |\nabla\Phi_0| > 0$
away from a centre, $\nuh = \nabla\Phi_0/g$ for the unit normal to the
level surfaces (towards increasing $\Phi_0$), and
$\nabla_{\!\parallel}$ for the part of the gradient along them. The
level set $\{\Phi_0 = c\}\cap F$ is $L_c$; the region below it,
$\{\Phi_0 < c\}\cap F$, is $R_c$; the part of the wall below it,
$\Sigma\cap\{\Phi_0 < c\}$, is $\Sigma_c$. The solid may be anything:
aspherical, pre-stressed, non-hydrostatic (\S\ref{sec:solid}); what it
must satisfy at the wall is stated there.

\paragraph{Level and transverse walls.} The wall is \emph{level}
where it is a level surface of $\Phi_0$ ($\bn = \nuh$), and
\emph{transverse} where it cuts the level surfaces. An enclosed fluid
is in equilibrium against a wall of any shape, the solid supplying the
traction; only a free surface must be level. On a spherical body, and
on any hydrostatic equilibrium figure, every density interface is
level. The Earth's core--mantle boundary is level to about a part in a
thousand.

\paragraph{Does the level foliation exist?} The closure needs
$\nabla\Phi_0 \ne \mathbf 0$ in $F$ away from a centre, and one
connected component per level set, so that $V(c)$ in
\S\ref{sec:closure} is a single function on an interval. Neither is
automatic. What is known:
\begin{itemize}
\item \emph{The maximum principle (derived).} Inside matter $\Phi_0$
  is subharmonic, $\nabla^2\Phi_0 = 4\pi G\rho > 0$ (with rotation
  $4\pi G\rho - 2\Omega^2$, still positive short of break-up), so it
  has no interior local maximum. Every component of a super-level set
  reaches the outer surface, and the foliation can degenerate only
  through interior minima (centres of attraction inside $F$) and the
  saddles between them. A single body near hydrostatic equilibrium
  has one minimum, at its centre, by continuity from the spherical
  case; a fluid shell then contains no critical point, and a full
  core exactly one.
\item \emph{The hydrostatic condition upgrades this (derived).} With
  $\rho = \rho(\Phi_0)$, $\Phi_0|_F$ solves the semilinear equation
  $\nabla^2\Phi_0 = 4\pi G\rho(\Phi_0)$ in $F$; on a level wall the
  boundary datum is the constant $c_w$, so $\Phi_0|_F$ is the solution
  of a Dirichlet problem on $F$ alone. With $w = c_w - \Phi_0 > 0$,
  $\nabla^2 w + 4\pi G\rho = 0$, $w = 0$ on $\Sigma$: a torsion-type
  problem, to which the convexity-of-level-sets theory applies.
\item \emph{A full core in a convex region (cited).} For uniform
  density this is the torsion problem, and $\sqrt w$ is concave
  (Makar-Limanov 1971 in two dimensions, Kennington 1985 in any): the
  level sets of $\Phi_0$ are nested convex surfaces with the single
  critical point. For stratified density, Kennington's Theorem~4.2
  (as restated by Ishige and Salani 2014) applies to
  $\nabla^2 w + f(\bx) = 0$ with $f = 4\pi G\rho(\bx)$ read as a
  function of position: if the density field is concave on $F$ (more
  generally $q$-concave for some $q \ge 1$), then $w^{q/(1+2q)}$ is
  concave, hence $w$ is quasi-concave and the conclusion is the same.
  A density profile that is a concave non-increasing function of
  radius on a ball --- uniform, linear in $r$, the stable cores of the
  benchmark ladder --- satisfies this; a profile convex in $r$ is not
  covered. On a ball the Gidas--Ni--Nirenberg theorem makes $\Phi_0$
  radial and monotone outright. Lions's conjecture that every
  non-negative nonlinearity on a convex domain gives convex level sets
  is false (Hamel, Nadirashvili and Sire 2016), so no blanket
  statement exists.
\item \emph{A shell (cited; no theorem).} An outer core between two
  level walls is a ring problem with constant data on each wall. The
  ring theorem of Caffarelli and Spruck 1982 requires, in the
  convention $\nabla^2 w = f(w)$, $f$ non-decreasing with $f(0) = 0$,
  which for the gravitational shell would mean density increasing
  outward and vanishing at the outer wall: not the physical case. And
  Hamel, Nadirashvili and Sire 2016 (Theorem 1.4) show that with $f$ a
  positive constant --- uniform density --- there are convex rings on
  which the level sets are not convex. So for the shell the assurance
  is near-sphericity, not a theorem; the counterexamples need rings
  far from concentric spheres.
\item \emph{Transverse walls.} The boundary datum is no longer
  constant and none of the above applies.
\end{itemize}
The implementation therefore asserts the assumption rather than
relying on it (\S\ref{sec:implementation}). If it fails, the
formalism survives on the Reeb graph of $\Phi_0|_F$ rather than on an
interval: one $V$ per component, $V$ additive and $f$ continuous
(pressure continuity) where components merge, and $A(c)$ with an
integrable singularity at a saddle level. That is bookkeeping, not
new theory, and is not built.

\paragraph{The fluid.} The fluid is inviscid and barotropic-elastic:
its strain energy depends on position and on the volume ratio only,
with bulk modulus $\kappa$ at equilibrium, and it supports no shear
stress. The position dependence carries the stratification (an
entropy or composition frozen into the material). The squared
buoyancy frequency is
\begin{equation}
  N^2 = -g^2\Bigl(\frac{\rho_\Phi}{\rho} + \frac{\rho}{\kappa}\Bigr),
  \qquad
  D := \rho_\Phi + \frac{\rho^2}{\kappa} = -\frac{\rho N^2}{g^2},
  \label{eq:N2}
\end{equation}
which on a radial background is the usual
$N^2 = -g(\rho'/\rho + g\rho/\kappa)$. $N^2 = 0$ is the neutral
(Adams--Williamson) state; the closure below never divides by $N^2$ or
$D$.

\paragraph{Perturbations.} $\bu$ is the displacement; $\phi$ the
Eulerian perturbation of the gravitational potential (the change at a
fixed spatial point), with $\nabla^2\phi = 4\pi G\rho^1$; in the fluid
$p^1$ and $\rho^1$ are the Eulerian pressure and density perturbations.
The Lagrangian pressure perturbation is $\Delta p = p^1 +
\bu\cdot\nabla p^0 = -\kappa\Div\bu$, and mass conservation is
$\rho^1 = -\Div(\rho\bu)$. The linearised static momentum balance of
the fluid is
\begin{equation}
  -\nabla p^1 - \rho\,\nabla\phi - \rho^1\,\nabla\Phi_0 = \mathbf 0 .
  \label{eq:momentum}
\end{equation}
The Lagrangian perturbation of the potential, the change of
$\Phi_0 + \phi$ following a parcel, is $\Delta\Phi = \phi +
\bu\cdot\nabla\Phi_0$.

\paragraph{The solid.} The solid's displacement is $\bu_s$. Two
descriptions of the solid are used: the \emph{mixed} one (referential
displacement, Eulerian potential; hydrostatic and natural reference
state; the class \code{LinearQuasiStaticMixedSelfGravitatingProblem}),
and the general \emph{referential} one (referential displacement and
potential; arbitrary constitutive state
$(\hat C, \Se, \varphi_e)$; the class
\code{LinearQuasiStaticReferentialSelfGravitatingProblem}). Sections
\ref{sec:dahlen}--\ref{sec:closure} are written in the mixed
description; \S\ref{sec:solid} lifts the restriction.

\paragraph{Level-surface integrals.} For a function $h$ on the fluid
region,
\begin{equation}
  A(c) = \oint_{L_c}\frac{\dS}{g},
  \qquad
  \bar h(c) = \oint_{L_c}\frac{h\dS}{g},
  \qquad
  \frac{\mathrm d}{\mathrm dc}\int_{R_c}h\dV = \bar h(c)
  \quad\text{(coarea)},
  \label{eq:coarea}
\end{equation}
so $A$ is the volume per unit potential and $\bar\Phi$, $\bar X$ are
the level means (weighted by $1/g$) of $\phi$ and of a test function
$\chi$. No level surface is ever constructed in the implementation:
every level integral is converted back to a volume integral by
\eqref{eq:coarea} (\S\ref{sec:implementation}).

%=====================================================================
\section{Why the fluid must be eliminated}
\label{sec:why}

\emph{Status: argued and measured (in \code{doc/static\_fluid\_core.tex}).}
The natural reading of ``static'' --- set the time derivatives to zero
and keep the fluid as an elastic material with $\mu = 0$, welded or
slipping at the wall --- has no solution when $N^2 \ne 0$. The
momentum balance alone forces the fluid's displacement along
$\nabla\Phi_0$ to be slaved to the potential, $\Delta\Phi = 0$, so
that fluid parcels ride the perturbed equipotentials
(\S\ref{sec:general}, with $f = 0$); the kinematic matching at the
wall, $\bu\cdot\bn = \bu_s\cdot\bn$, asks the same trace to take a
second, generically different value. The discrete problems are a
minimising sequence whose infimum is not attained: under refinement the
computed fluid displacement develops ever finer structure at the wall
and does not converge in the displacement space, while the solid
observables stay within a band set by the stratification, which no
resolution removes. Dissipation of any kind
(viscosity, diffusion of the stratifying quantity, a Maxwell core)
gives a well-posed evolution whose long-time limit is the state in
which the fluid has rearranged to hydrostatic equilibrium in the
perturbed potential. That \emph{secular state} is what a static fluid
method must compute, and it is what Dahlen's elimination computes:
the fluid's displacement is never introduced, so the conflict at the
wall never arises, and the mass the slaved field would carry across
the wall is booked as a surface density. The closure of
\S\ref{sec:closure} keeps this structure on any geometry.

%=====================================================================
\section{Dahlen's reduction on a spherical body}
\label{sec:dahlen}

\emph{Status: derived (the reduction); measured (its properties).}
Let the background be radial, so that the level surfaces are the
spheres, $\nabla\Phi_0 = g\hat{\mathbf r}$, and every perturbation
separates into spherical harmonics (circular harmonics in two
dimensions) of degree $\ell$.

\paragraph{The interior relations.} The tangential part of
\eqref{eq:momentum} has no $\rho^1$ term, since $\nabla\Phi_0$ is
radial, and $\rho$ is constant on each sphere, so
\begin{equation}
  \nabla_{\!\parallel}\bigl(p^1 + \rho\phi\bigr) = \mathbf 0,
  \qquad
  p^1 + \rho\,\phi = f(r) .
  \label{eq:tangential}
\end{equation}
A function of $r$ alone has no component along a harmonic of degree
$\ell \ge 1$, so at every such degree
\begin{equation}
  p^1 = -\rho\,\phi \qquad (\ell \ge 1).
  \label{eq:p-slaved}
\end{equation}
The radial part of \eqref{eq:momentum} then gives, with
$\partial_r p^1 = -\rho'\phi - \rho\,\partial_r\phi$,
\begin{equation}
  \rho^1 = \frac{\rho'}{g}\,\phi = \rho_\Phi\,\phi \qquad (\ell \ge 1).
  \label{eq:rho-slaved}
\end{equation}
These two relations follow from momentum balance alone; they hold for
any static description of the fluid. Dahlen's reduction keeps them,
discards the constitutive law and the kinematics, and never introduces
a fluid displacement.

\paragraph{The mixed weak form.} The solid is described in the mixed
organisation; the fluid contributes three terms, named F1--F3 in this
library (\code{mixed\_problem.hpp}; \code{doc/benchmarks.tex}, ``The
CMB treatments''). With $\mathbf m = -\bn$ the solid's outward normal
on $\Sigma$ and $\rho_F$ the fluid-side density there,
\begin{align}
  &\text{(F1)}\quad \int_F\rho_\Phi\,\phi\,\chi\dV
  &&\text{the fluid's Eulerian density perturbation
     \eqref{eq:rho-slaved} in Poisson's equation,}
  \nonumber\\
  &\text{(F2)}\quad -\int_\Sigma\rho_F\,(\mathbf m\cdot\nabla\Phi_0)
    (\mathbf m\cdot\bu_s)(\mathbf m\cdot\bv)\dS
  &&\text{the buoyancy of the wall: } \bu_s\cdot\nabla p^0
    \text{ in the pressure,}
  \label{eq:F123}\\
  &\text{(F3)}\quad -\int_\Sigma\rho_F\,\bigl[\phi\,(\mathbf m\cdot\bv)
    + \chi\,(\mathbf m\cdot\bu_s)\bigr]\dS
  &&\text{the pressure } -\rho_F\phi \text{ on the solid, and the surface
     mass } \rho_F\,\bu_s\cdot\bn .
  \nonumber
\end{align}
The full weak form of the class is: the solid's elastic and
hydrostatic-gravity terms plus F2 in the displacement row; the
Laplacian with its exterior condition plus F1 in the potential row; and
the coupling $c(\phi, \bv) = \int_{\text{solid}}\rho\nabla\phi\cdot\bv -
\int_\Sigma\rho_F\phi(\mathbf m\cdot\bv)$, whose interface part is F3,
in both. The fluid's displacement and its bulk modulus appear
nowhere.

\paragraph{What the elimination hides.} Reconstruct the displacement
implied by \eqref{eq:rho-slaved} through $\rho^1 = -\Div(\rho\bu)$: it
is the slaved field $u_r = -\phi/g$, that is $\Delta\Phi = 0$, which
does not match the solid's $u_r$ at the wall. The formalism is
consistent because of how F3 books the interface. Its surface density
is $\rho_F\,\bu_s\cdot\bn$, the fluid-side half of the density-jump
term $(\rho_F - \rho_S)\,\bu_s\cdot\bn$ that a displacement continuous
across the wall would produce (the solid-side half, $-\rho_S\bu_s\cdot\bn$,
sits inside the solid's volume coupling after integration by parts).
It is not the surface density the slaved fluid itself would produce,
which is $\rho_F u_\nu$ with $u_\nu = -\phi/g$, the fluid's own boundary
motion. The difference, $\rho_F(\bu_s\cdot\bn - u_\nu)$, is mass added
where the wall has moved away from the slaved fluid and removed where
it has moved into it: the normal-displacement discontinuity, booked as
a surface mass. So F3 does more than move the boundary; it fills the
gap, and the fluid \emph{fills to the wall}. In the displacement
picture this is the wall-level sheet of fluid spreading tangentially
along the wall into the gap; its Eulerian effect is exactly the
booking above, and its net mass is zero, since it is a redistribution.
That is the secular closure: the state after the fluid has rearranged,
however long it takes.

\paragraph{Properties (measured).} (i) At $\ell \ge 1$ the construction
never consults $\kappa$: Dahlen's answer is independent of the fluid's
bulk modulus (unchanged to $10^{-11}$ under a thirty-fold scaling).
(ii) At $N^2 = 0$ the reduction coincides with the full elastic
static fluid, so Dahlen, welded and slipping agree there
($\sim 10^{-10}$ at every degree on the neutral core). (iii) Where
$N^2 \ne 0$ Dahlen converges to machine precision while the
displacement formulations wander in their band.

\paragraph{The degree-0 hole.} At $\ell = 0$ the argument of
\eqref{eq:p-slaved} fails at its first line: every degree-0 field is a
function of $r$, the tangential balance is empty, and $f(r)$ in
\eqref{eq:tangential} is a free function. The reduction then leaves
one function undetermined, and a discretisation of F1--F3 alone
computes an arbitrary member of the family. Measured: on the
\code{fluid\_core} model the mixed class gives $h'_0 = -1.457$ and
$k'_0 = -5.35$ against the reference $-0.1321$ and $0$ (mass
conservation fixes $k'_0 = 0$); on the two-dimensional disc the
reduction alone gives $u_r(a) = +3.29$ against $-0.0991$. The
spectral codes avoid the hole by solving $\ell = 0$ with the
displacement description, where the harmonics decouple; a
three-dimensional discretisation has no such sector. Closing the hole
within the three-dimensional problem is one of the two things the
generalisation does.

%=====================================================================
\section{A hydrostatic fluid region of any shape}
\label{sec:general}

\emph{Status: derived.} The spherical reduction used the sphere
twice: level surfaces were coordinate spheres, and harmonics
decoupled. Neither is needed for the interior relations.

\paragraph{The free function.} The along-level part of
\eqref{eq:momentum} has no $\rho^1$ term ($\nabla_{\!\parallel}\Phi_0 =
\mathbf 0$) and $\rho$ is constant on each level component, so
\begin{equation}
  p^1 + \rho\,\phi = f(\Phi_0),
  \label{eq:G-f}
\end{equation}
with $f$ a free function, one value per connected component of each
level set, continuous where components merge. The normal part, with
$\partial_\nu\rho = \rho_\Phi g$ and $\partial_\nu f = f' g$ (primes
are $\mathrm d/\mathrm d\Phi_0$ from here on), gives
\begin{equation}
  \rho^1 = \rho_\Phi\,\phi - f'(\Phi_0) .
  \label{eq:G-rho}
\end{equation}
On the sphere at $\ell \ge 1$, $f = 0$ and these are
\eqref{eq:p-slaved}--\eqref{eq:rho-slaved}. Dahlen's reduction is
\eqref{eq:G-f}--\eqref{eq:G-rho} with $f$ left free; the general
problem is to determine $f$.

\paragraph{The kinematics, if a displacement is added.} Define the
\emph{lift} $s$ by $\bu\cdot\nabla\Phi_0 = -\phi - s$, so $s =
-\Delta\Phi$. Then $\Delta p = f + \rho s$, $\Div\bu = -(f + \rho s)/\kappa$,
and equating the two determinations of $\rho^1$ (mass conservation and
\eqref{eq:G-rho}),
\begin{equation}
  D\,s = -\Bigl(f' + \frac{\rho f}{\kappa}\Bigr),
  \qquad
  \Div\bu = -\frac{f + \rho s}{\kappa} .
  \label{eq:G-slaving}
\end{equation}
The normal displacement and the dilatation are slaved pointwise to
$f$ and $f'$. If $\kappa$ is constant on the level surfaces, so are
$s$ and $\Div\bu$: material surfaces that were level stay level in the
perturbed potential. (With $f = 0$ this is the slaving of
\S\ref{sec:why}, $s = 0$.) $\kappa$ \emph{is} constant on the level
surfaces for any real fluid: with $\rho = \rho(p, S)$ and
$\partial_S\rho \ne 0$, level-constant $\rho$ and $p^0$ force
level-constant $S$, hence $\kappa$. The along-level displacement
$\bu_\parallel$ is not slaved; it is determined, up to the
relabellings that are tangent to the level surfaces and
divergence-free, by its surface divergence on each level surface, and
that divergence is fixed by $\Div\bu$ and the normal component.

\paragraph{The one condition per level.} A tangential field on a
level surface with prescribed surface divergence and prescribed flux
through the surface's boundary (on the wall) exists if and only if the
two are compatible: one scalar condition per level component. That
condition is mass conservation of the region below the level, with the
no-penetration flux on the wall patch,
\begin{equation}
  \delta M(c) = \int_{R_c}\rho^1\dV
  + \rho(c)\,U(c)
  + \int_{\Sigma_c}\rho_F\,\bu_s\cdot\bn\dS = 0,
  \qquad
  U(c) = \oint_{L_c}u_\nu\dS = -\bigl(\bar\Phi + sA\bigr),
  \label{eq:G-band}
\end{equation}
for every $c$. There are as many conditions as values of $f$, so the
over-determination of \S\ref{sec:why} does not occur once $f$ is free:
the interior slaving and the wall can be reconciled by the right $f$.
On a level wall $\Sigma_c$ is empty below the wall's level, and the
wall itself supplies one further condition (\S\ref{sec:closure}).

%=====================================================================
\section{The closure in the cumulative-volume variable}
\label{sec:closure}

\emph{Status: derived; measured in one dimension.} The free function is
best carried not as $f$ but as the \emph{cumulative volume change below
each level},
\begin{equation}
  V(c) = \int_{R_c}\Div\bu\dV = U(c) + S_c[\bu_s],
  \qquad
  S_c[\bu_s] = \int_{\Sigma_c}\bu_s\cdot\bn\dS ,
  \label{eq:V-def}
\end{equation}
by the divergence theorem with no-penetration on the wall patch. By
the coarea formula and the level-constancy of $\Div\bu$,
$V' = A\,\Div\bu$. Then from \eqref{eq:G-band}, \eqref{eq:G-slaving}
and \eqref{eq:G-rho},
\begin{equation}
  s = -\frac{\bar\Phi + V - S_c[\bu_s]}{A},
  \qquad
  f = -\frac{\kappa V'}{A} - \rho\,s,
  \qquad
  \rho^1 = \rho_\Phi(\phi + s) - \frac{\rho V'}{A},
  \label{eq:V-dict}
\end{equation}
and the first relation of \eqref{eq:G-slaving} is an identity. The one
equation left is the normal balance $\rho^1 = \rho_\Phi\phi - f'$, that
is $f' = \rho V'/A - \rho_\Phi s$. On a level wall ($S_c = 0$ below the
wall's level $c_w$) it reads
\begin{equation}
  \boxed{\;-\Bigl(\frac{\kappa V'}{A}\Bigr)' + \rho\Bigl(\frac1A\Bigr)'V
  + \rho\Bigl(\frac{\bar\Phi}{A}\Bigr)' = 0\;}
  \qquad c_{\min} < c < c_w,\qquad V(c_{\min}) = 0 ,
  \label{eq:V-ode}
\end{equation}
a second-order equation in $V$ driven by the level means of $\phi$.
(On a shell between two level walls the same equation holds for the
flux $U$ --- $V$ and $U$ differ there by the constant volume swept by
the inner wall, which drops out of every derivative --- with a wall
constraint at each end; the implementation carries $U$ for this
reason.) Nothing divides by $D$ or $N^2$: the closure is regular where
the stratification changes sign and at $N^2 = 0$. Two checks on the
sphere (derived): with $V = 4\pi r^2u_r$, $A = 4\pi r^2/g$ and
$\mathrm d/\mathrm dc = g^{-1}\mathrm d/\mathrm dr$, the dead-gravity
limit $\phi = 0$ of \eqref{eq:V-ode} is $(\kappa\Div\bu)_r +
\rho u_r(2g/r - g') = 0$, the degree-0 radial equation of the
$\mu = 0$ fluid; and \eqref{eq:V-dict} gives $f = p^1 + \rho\phi$.

\subsection{The wall condition}
\label{sec:wall}

On a level wall the fluid's slaved normal trace $u_\nu = -(\phi + s)/g$
and the solid's $\bu_s\cdot\bn$ differ pointwise; their difference is
the surface density of the spread sheet of \S\ref{sec:dahlen},
$\sigma = \rho_F(u_\nu - \bu_s\cdot\bn)$. The sheet redistributes
wall-level fluid along the wall and carries no net mass, since every
band below conserves its own mass by \eqref{eq:G-band}. Hence
\begin{equation}
  \oint_\Sigma\sigma\dS = 0
  \qquad\Longleftrightarrow\qquad
  V(c_w) = S[\bu_s] := \oint_\Sigma\bu_s\cdot\bn\dS ,
  \label{eq:wall}
\end{equation}
the \emph{level mean of the normal matching}: the fluid's volume flux
through the wall level equals the volume swept by the wall. It is
\eqref{eq:G-band} applied to the whole region, $\int_F\rho^1 +
\oint\rho_F\bu_s\cdot\bn = 0$, the fluid mass inside the displaced
container, surface mass excluded, being conserved. Structurally it is
an essential constraint whose Lagrange multiplier is the level-mean
wall pressure (\S\ref{sec:weak}); without it the solid's wall-pressure
column has no transpose and the system is not square. On the sphere at
$\ell = 0$, $\sigma$ is uniform, so \eqref{eq:wall} is $\sigma = 0$:
continuity of $u_r$, the degree-0 matching.

\subsection{The bordered weak form}
\label{sec:weak}

Write $A_w = A(c_w)$, $\bar X(c)$ for the level mean of a potential
test function $\chi$, and $L_D(\bu_s, \phi)$ for the quadratic
functional of the mixed description with F1--F3 included. The closure
is the stationarity of
\begin{equation}
\begin{aligned}
  L = L_D(\bu_s,\phi)
  &+ \frac12\int\frac{\kappa V'^2}{A}\dc
   + \frac12\int\rho\Bigl(\frac1A\Bigr)'V^2\dc
   - \int\Bigl(\frac{\rho_\Phi\bar\Phi V}{A}
            + \frac{\rho\,\bar\Phi V'}{A}\Bigr)\dc
   - \frac12\int\frac{\rho_\Phi\bar\Phi^2}{A}\dc \\
  &- \frac{\rho_F}{2A_w}\,S[\bu_s]^2
   + \lambda\bigl(V(c_w) - S[\bu_s]\bigr) - \text{loads},
\end{aligned}
\label{eq:L}
\end{equation}
with $V(c_{\min}) = 0$ essential and $\lambda$ a scalar multiplier.
With test functions $\bv$, $\chi$, $W$ ($W(c_{\min}) = 0$) the rows
are
\begin{equation}
\begin{aligned}
  \text{(S)}&\quad
  \text{mixed solid row} - \lambda\,S[\bv]
  - \frac{\rho_F}{A_w}S[\bu_s]\,S[\bv] = \ell_u(\bv),\\
  \text{(P)}&\quad
  \text{mixed potential row}
  - \int\frac{\rho_\Phi\bar\Phi\bar X}{A}\dc
  - \int\Bigl(\frac{\rho_\Phi V}{A} + \frac{\rho V'}{A}\Bigr)\bar X\dc
  = \ell_\phi(\chi),\\
  \text{(V)}&\quad
  \int\frac{\kappa V'W'}{A}\dc + \int\rho\Bigl(\frac1A\Bigr)'VW\dc
  - \int\Bigl(\frac{\rho_\Phi\bar\Phi W}{A} + \frac{\rho\,\bar\Phi W'}{A}\Bigr)\dc
  + \lambda\,W(c_w) = 0,\\
  \text{(C)}&\quad V(c_w) - S[\bu_s] = 0 .
\end{aligned}
\label{eq:rows}
\end{equation}
The system is symmetric by construction: the (V)--(P) coupling is one
bilinear form read both ways, and the (V) row is the weak form of
\eqref{eq:V-ode} with its natural flux absorbed into the multiplier,
\begin{equation}
  \lambda = -\frac{\kappa V'(c_w)}{A_w} + \frac{\rho_F\bar\Phi(c_w)}{A_w}
  = f(c_w) - \frac{\rho_F S[\bu_s]}{A_w} .
  \label{eq:lambda}
\end{equation}
So (S) carries exactly the wall pressure $f(c_w)\oint(\mathbf m\cdot\bv)$
once the rank-one term is accounted for: the solid feels the fluid
through $f - \rho_F(\phi + \bu_s\cdot\nabla\Phi_0)$ at the displaced
wall, the last two pieces being F3 and F2.

\paragraph{Bookkeeping on the sphere.} At $\ell = 0$ the term
$-\tfrac12\int\rho_\Phi\bar\Phi^2/A$ cancels the degree-0 part of F1
exactly, and $-(\rho_F/2A_w)S^2$ cancels the degree-0 part of F2: both
are consequences of the slaving $u_r = -\phi/g$, which does not hold
at degree 0. F3 stays, being the boundary term that the welded
coupling $\int_F\rho\nabla\phi\cdot\bu$ produces under integration by
parts. So at $\ell = 0$ the bordered system is the welded compressible
fluid in a different assembly, and at $\ell \ge 1$ the border rows are
homogeneous ($\bar\Phi = 0$, $S[\bu_s] = 0$), $V = \lambda = 0$, and the
method is Dahlen's unchanged. Off the sphere the level means of a
general load do not vanish, and the border is active for every load.

\subsection{What is settled, and what is not}
\label{sec:settled}

\begin{itemize}
\item \textbf{Uniqueness} is the coercivity of \eqref{eq:L} on the
  complement of the rigid modes: the gravitational stability of the
  level-mean sector. The sign of $N^2$ does not enter (there is no
  overturning in that sector), and at $N^2 = 0$ the sector carries no
  relabelling freedom, since a level-mean relabelling has
  $w_r \propto 1/(\rho r^2)$, singular at a centre and annihilated by
  the walls of a shell. On the sphere it is the coercivity of the
  degree-0 problem. Off the sphere it is \emph{argued} by continuity
  and measured by the solver's minimum-eigenvalue probe.
\item \textbf{The along-level displacement} is never constructed. It
  carries no energy and enters no Eulerian field beyond its divergence,
  which \eqref{eq:V-def} fixes. Its growth like $1/\sin\theta$ near a
  level point of a transverse wall ($\theta$ the angle between wall and
  level surface) is immaterial to the reduced problem.
\item \textbf{F3 needs no limiting argument.} It books the fluid
  filling to the wall, on level and transverse walls alike.
\item \textbf{The welded count} on a transverse wall is not a question:
  an inviscid fluid slips, and the secular state has a free tangential
  trace.
\end{itemize}

\subsection{Transverse walls}
\label{sec:transverse}

\emph{Status: argued in outline.} Nothing changes in form. The strips
$\Sigma\cap\{\Phi_0 = c\}$ enter through $U = V - S_c[\bu_s]$ in
\eqref{eq:V-dict}, the wall pressure $f(\Phi_0(\bx))$ acts along the
wall, the wall's normal displacement enters the (V) row as
$\int_\Sigma W(\Phi_0)\,\rho_F\,\bu_s\cdot\bn\dS$, and
$V(c_{\mathrm{top}}) = S[\bu_s]$ remains the end condition; the level
wall is the case in which the strips concentrate at $c_w$. The
transverse case is regular in $V$. But a density jump across a
transverse wall is incompatible with a hydrostatic solid: traction
continuity along the wall requires $p^0_S(c) = p^0_F(c)$ over the cut
range of $\Phi_0$, whose $c$-derivative gives $\rho_S = \rho_F$ there.
A transverse wall with a density contrast needs deviatoric pre-stress
in the solid, so it belongs to \S\ref{sec:solid}, not to the mixed
description. One prediction for the displacement formulations follows:
on a genuinely transverse wall the elastic-static slipping problem has
a regular solution except near the level points every closed wall has,
so its resolution band should shrink with the wall's transversality.

%=====================================================================
\section{A general pre-stressed solid}
\label{sec:solid}

\emph{Status: argued.} Nothing in \S\ref{sec:general}--\ref{sec:closure}
consults the solid. The fluid relations are Eulerian and local to $F$,
and the solid enters them through two quantities only: the total
background potential $\Phi_0$, sourced by every mass including a
pre-stressed solid, and the wall's normal displacement. The closure
therefore couples unchanged to the general referential equations in
$(\hat C, \Se, \varphi_e)$ (\code{doc/gravitating\_elasticity.md}
\S3.1), and the mixed description is the special value
$\Se = -p^0 I$, $\varphi_e = \mathrm{id}$. The organisation is hybrid:
the referential potential $\zeta^1$ on the solid and the vacuum, the
Eulerian potential on the fluid. It is also the natural organisation
for secular motions on an aspherical reference, since the referential
solid carries no bare displacement in its gravity coupling and the
equipotential-riding fluid carries no displacement at all.

\paragraph{The potential across the wall.} On the solid
$\zeta^1 = \phi^1\circ\varphi_e + (\bu\cdot\nabla\Phi_0)\circ\varphi_e$;
in the fluid, with no $\bu$, the same formula gives
$\phi^1\circ\varphi_e$. Two realisations exist. Either the displacement
is extended into the fluid by the prescribed radial extension $E$ of
the slipping classes and every fluid term is written in
$\phi^1 = \zeta^1 - (E\bu)\cdot\nabla\Phi_0$, an exact change of
variables with no gauge content since $E\bu$ is never treated as a
fluid displacement; or $\zeta^1$ is broken at $\Sigma$ with the jump
prescribed as $\bu_s\cdot\nabla\Phi_0$. The first keeps one potential
space, and the referential Poisson operator on the fluid region is
already the mapped Laplacian the class assembles.

\paragraph{The fluid terms, pulled back.} The volume source and the
border integrals are assembled on the reference fluid mesh with the
basis functions composed with $\Phi_0\circ\varphi_e$ and the Jacobian
$J_e$; the wall integrals carry the surface Jacobian; the constraint
reads the referential flux through the cofactor normal.

\paragraph{The follower load.} This is the one term the mixed
description never shows, because a hydrostatic solid absorbs it into
its symmetrised advective form. The fluid's traction on the deformed
wall, $-p\,\bn\dS$ linearised, splits into the pressure perturbation
at the displaced wall, $\delta p = f - \rho_F(\phi^1 + \bu_s\cdot\nabla\Phi_0)$
on the reference normal, and $-p^0$ times the perturbation of the
normal-area element. The second is the one-sided solid kernel of the
slipping classes' interface pressure form at $\pi = p^0$
(\code{SlipInterfacePressureIntegrator}). On a level wall $p^0$ is
constant along $\Sigma$ and that kernel is symmetric by itself. On a
transverse wall it is not, and the total is symmetric only when
\begin{equation}
  \Se\,\bn = -p^0\,\bn \quad\text{on } \Sigma :
  \label{eq:traction}
\end{equation}
the solid's pre-stress carries the fluid's normal traction and no shear
traction on the wall. The equilibrium-stress generators of
\code{background.hpp} take boundary tractions as data, so admissible
$\Se$ for a non-level wall are constructible, the wall pressure being
the datum.

\paragraph{The fluid half of F3.} The solid's $-\rho_S\,\bu_s\cdot\bn$
is inside the referential potential row automatically ($\rho$ is a
fixed referential field); the fluid's $+\rho_F\,\bu_s\cdot\bn$ is added
explicitly, as in the mixed description.

\paragraph{Consistency conditions.} (i) $\Phi_0$ is the total potential
of the equilibrium configuration, the pre-stressed solid included.
(ii) The fluid is hydrostatic in it. (iii) \eqref{eq:traction} on
$\Sigma$. (iv) The solid is otherwise arbitrary. Under (i)--(iii) the
referential solid with the Eulerian fluid closure and its border is a
symmetric bordered system with the variational structure of the
conservative system solid, fluid and gravity. Not yet written out term
by term: the follower term's symmetry with a deviatoric $\Se$ on a
transverse wall, and the pull-back of the border under
$\varphi_e \ne \mathrm{id}$.

%=====================================================================
\section{Rotation}
\label{sec:rotation}

\emph{Status: derived for the fluid; argued for the meeting with the
generalised rotational theory.}

\paragraph{The background.} No fluid relation above used Poisson's
equation for $\Phi_0$, only \eqref{eq:hydro} and Poisson's equation for
the perturbation $\phi$. So with $\Phi_0$ the \emph{geopotential}
(gravitational plus centrifugal, $-\tfrac12|\bOmega\times\bx|^2$),
$g = |\nabla\Phi_0|$ the effective gravity and $N^2$ formed with it,
every step holds verbatim: the fluid is foliated by the geopotential
surfaces. On a rotating hydrostatic figure the body is axisymmetric
about $\bOmega$ and the core--mantle boundary is a geopotential
surface, the level-wall case.

\paragraph{The perturbation.} A change $\bomega$ of the rotation vector
perturbs the centrifugal potential by
$\psi = -(\bOmega\times\bx)\cdot(\bomega\times\bx) = \sum_k\omega_k\psi_k$,
given once $\bomega$ is, and $\psi$ joins $\phi$ everywhere in the
fluid: $f = p^1 + \rho(\phi + \psi)$, $\rho^1 = \rho_\Phi(\phi + \psi) - f'$,
$s = -(\phi + \psi + \bu\cdot\nabla\Phi_0)$, the level means
$\bar\Phi \to \bar\Phi + \sum_k\omega_k\bar\Psi_k$, and the wall pressure
$f - \rho_F(\phi + \psi + \bu_s\cdot\nabla\Phi_0)$. Making that
substitution in \eqref{eq:L} gives the border its $\bomega$ columns,
$-\int(\rho_\Phi\bar\Psi_kV/A + \rho\bar\Psi_kV'/A)\dc$ in the (V) row
and $-\int\rho_\Phi\bar\Phi\bar\Psi_k/A\dc$ in the (P) row, and an
$\bomega$ block $-\tfrac12\int\rho_\Phi\bar\Psi_j\bar\Psi_k/A\dc$; by
the symmetry of the functional the $\bomega$ rows acquire the
transposes. This is the route by which the tidal machinery gives the
mixed description its $\bomega$-row data: that row is the fluid's
inertia perturbation $\int_F\rho^1(|\bx|^2\mathbf 1 - \bx\otimes\bx)\dV$
plus the surface-mass term, paired with $\bOmega$, and the border
supplies the $-f'$ part of $\rho^1$ that the reduction alone lacks.

\paragraph{A consequence for the reduction alone (derived).} A change
of rotation rate, $\bomega \parallel \bOmega$, has
$\psi = -\Omega\omega_3(x^2 + y^2)$, whose degree-0 part on the sphere
is $-\tfrac23\Omega\omega_3r^2$. The $\omega_3$ row of the secular
balance, $-C_3\omega_3 = \Omega I^{(1)}_{33}$, reads the degree-0 part
of $\rho^1$ through $I^{(1)}_{33} = \int\rho^1(x^2 + y^2)$, and the
feedback $\psi(\omega_3)$ loads degree 0 in turn. With the reduction
alone the fluid's degree-0 response is the arbitrary family member of
\S\ref{sec:dahlen}, so the length-of-day response of a fluid-core
model under a load with a degree-0 part is contaminated at the level of
that member; polar wander ($\bomega \perp \bOmega$) is degree 2 only and
is clean. The closure removes the contamination. This is checkable
against the spectral reference with no new theory.

\paragraph{The generalised rotational theory
(\code{doc/generalised\_rotation.md}).} The secular reduction of
Maitra and Al-Attar's theory keeps three things: the inertia matrix
$D$ built from the model's own equilibrium inertia $I^{(0)}$; the
displacement pinned by the Tisserand gauge; $\bOmega$ along a
principal axis. With the closure, all three survive unchanged.
$I^{(0)}$ is whole-body, the fluid's hydrostatic density included. The
Tisserand gauge restricted to the solid is forced rather than chosen:
the eliminated fluid has no displacement, and the slaved one's relative
angular momentum is relabelling gauge. $\bOmega$ along a principal axis
is automatic on an axisymmetric figure. The inertia perturbation is
Eulerian in the fluid with the border and surface-mass terms added.
Two things want verification rather than argument. First, the
gauge-shift certificate of that note must be re-derived: under a rigid
re-gauge $\bu_s \to \bu_s + \mathbf b\times\bx$ on an aspherical level
wall the fluid is not inert. Without rotation the Eulerian fluid
follows a rigid motion of the total mass exactly ---
$\rho^1 = -(\mathbf b\times\bx)\cdot\nabla\rho$ from $\rho_\Phi\phi$
with $\phi = -(\mathbf b\times\bx)\cdot\nabla\Phi_0$, and
$u_\nu = -\phi/g = (\mathbf b\times\bx)\cdot\bn$ on a level wall, so no
surface mass --- which is why a rigid rotation is a null pair of the
closure; with the centrifugal background it is a frame unknown, and
the expectation is that the whole-body commutator formula survives.
Second, the closure and the self-consistent door are made for each
other: a rotating hydrostatic figure with a fluid core is where the
model's own inertia is realistic and where, at $N^2 \ne 0$, no static
fluid method existed. The rapid-wobble regime, in which the core does
not follow the mantle, remains outside the secular door as before.

%=====================================================================
\section{The nonlinear theory}
\label{sec:nonlinear}

\emph{Status: the state space, the first variation and the dictionary
to the linear closure are derived; the identification of the reduced
second variation is argued; existence is cited; the rest is a
programme.} The closure of \S\ref{sec:closure} is the linearisation of
an exact theory in which the fluid has no displacement at any order.
The exact theory is the cleaner statement, it explains two features
of the linear closure that look accidental there, and it is the form
in which the method should eventually be built.

\subsection{The state space of the fluid}
\label{sec:NL-state}

Each fluid parcel carries a conserved scalar label $S$: specific
entropy for a thermally stratified fluid, composition for a chemically
stratified one, or any monotone function of either. The equation of
state is $\rho = \rho(p, S)$, with internal energy $e(\rho, S)$ per
unit mass, $p = \rho^2\partial_\rho e$, and bulk modulus
$\kappa = \rho\,\partial_\rho p|_S$. A configuration of the fluid is a
pair of fields $(\rho, S)$ on the region $\Omega_F$ bounded by the
deformed wall. The relabelling symmetry, exactly, is the group of
mass-preserving maps of the fluid that preserve the label; what it
leaves invariant is not a configuration but the \emph{mass--label
measure}
\begin{equation}
  \mu(E) = \int_{\{S \in E\}}\rho\dV ,
  \qquad
  m(S) = \mu\bigl((-\infty, S]\bigr),
  \label{eq:NL-mu}
\end{equation}
the mass of fluid carrying labels in $E$. Two configurations are
rearrangements of each other exactly when they have the same $\mu$,
and the physical state space of the fluid, the quotient by the
symmetry, is the \emph{rearrangement class} of $\mu$: all
$(\rho, S)$ on $\Omega_F$ with the given mass--label measure. This is
Burton's setting for steady flows as energy extremisers over
rearrangements, with the mass measure in place of the volume measure
because the fluid is compressible. In the Lagrangian picture the same
statement reads: a motion $\psi$ of a reference fluid body with
reference labels $S_0$ and density $\rho_0$ enters the energy only
through $(\rho_0/J)\circ\psi^{-1}$ and $S_0\circ\psi^{-1}$, so two
motions differing by a mass- and label-preserving map of the reference
body are the same state. The linear classes of
\code{doc/static\_fluid\_core.tex} --- the strict class at $N^2 \ne 0$,
class K at $N^2 = 0$ --- are the tangent spaces of this group at a
configuration: when $S$ is a non-trivial function the admissible
infinitesimal relabellings must preserve both $\rho$ and $S$, hence are
tangent to their common level surfaces and divergence-free; when $S$ is
constant (a barotropic fluid, $N^2 = 0$) only $\rho$ constrains them.

\subsection{The energy and the exact secular state}
\label{sec:NL-energy}

With $\varphi$ the solid's finite motion, the total energy is
\begin{equation}
  E[\varphi, \rho, S]
  = \int_{B_S}W(X, D\varphi)\dV
  + \int_{\Omega_F}\rho\,e(\rho, S)\dV
  - \frac{1}{8\pi G}\int_{\R^3}|\nabla\Phi_g|^2\dV
  + \int_{\Omega}\rho_{\mathrm{tot}}\,\Phi_c\dV ,
  \label{eq:NL-E}
\end{equation}
the solid's stored energy, the fluid's internal energy, the
gravitational self-energy of the total density (solid pushed forward,
fluid as it stands) with $\Phi_g$ its potential, and, for a body held
at angular velocity $\bOmega$, the centrifugal term with
$\Phi_c = -\tfrac12|\bOmega\times\bx|^2$; $\Phi = \Phi_g + \Phi_c$ is
the geopotential. (With $\bOmega$ dynamical the last term is replaced
by the rotational Routhian of Maitra and Al-Attar; nothing below
changes in the fluid.) The variables are the solid's motion and the
fluid's state in its rearrangement class, with the fluid filling the
container, $\Omega_F = \varphi(B_F)$-shaped by the wall; this is
no-penetration at finite amplitude.

\emph{First variation in the fluid (derived).} A variation within the
class is generated by a virtual displacement $\boldsymbol\xi$ of the
parcels, $\delta\rho = -\Div(\rho\boldsymbol\xi)$, $\delta S =
-\boldsymbol\xi\cdot\nabla S$, with $\boldsymbol\xi\cdot\bn$ equal to
the wall's normal displacement. The stationarity of \eqref{eq:NL-E}
under every such $\boldsymbol\xi$ is the hydrostatic balance
$\nabla p = -\rho\nabla\Phi$ in $\Omega_F$ and the traction $-p\,\bn$
on the solid's wall. Taking the curl, $\nabla\rho\times\nabla\Phi =
\mathbf 0$, and since $\rho = \rho(p, S)$ with $\partial_S\rho \ne 0$,
also $\nabla S\times\nabla\Phi = \mathbf 0$: in equilibrium the fluid's
state is a function of the level,
\begin{equation}
  S = \hat S(\Phi),\qquad p = \hat p(\Phi),\qquad
  \frac{\mathrm d\hat p}{\mathrm d\Phi} = -\rho\bigl(\hat p, \hat S\bigr),
  \label{eq:NL-profile}
\end{equation}
one monotone-or-not map and one pressure profile per connected
component of the level sets. The fluid carries the level foliation of
the geopotential of the deformed body, and nothing else.

\emph{The rearrangement constraint becomes one-dimensional (derived).}
What ties the profile \eqref{eq:NL-profile} to the fluid's identity is
that its state must lie in the class of $\mu$. Write $M(c) =
\int_{\{\Phi < c\}\cap\Omega_F}\rho\dV$ for the mass below level $c$;
by the coarea formula $M'(c) = \rho(c)\,A(c)$ with $A$ the level-area
density of \S\ref{sec:setting}. If $\hat S$ is monotone increasing, the
set of parcels with label below $\hat S(c)$ is the set below level
$c$, so
\begin{equation}
  M(c) = m\bigl(\hat S(c)\bigr)\quad\text{for every } c ,
  \label{eq:NL-mass}
\end{equation}
one scalar condition per level, and conversely \eqref{eq:NL-mass}
with \eqref{eq:NL-profile} determines $\hat S$ and $\hat p$ from the
level geometry $A(c)$: two first-order equations,
$\hat p' = -\rho(\hat p, \hat S)$ and $m'(\hat S)\hat S' = \rho(\hat p, \hat S)A$,
on the level range of the region. (For a monotone decreasing $\hat S$
the same holds with the order of labels reversed; the non-monotone
case is \S\ref{sec:NL-finite}.) The conditions are two. At the
bottom, a centre or an inner wall, $M = 0$ fixes $\hat S$ at the
lowest label. At the top, the fluid must fit its container:
$M(c_{\mathrm{top}}) = m(S_{\max})$, the total mass, which fixes the
one constant the pressure profile carries. That constant is the
degree-0 freedom of \S\ref{sec:dahlen}, and the condition that fixes
it is the nonlinear form of the wall condition \eqref{eq:wall}. The
exact secular state is therefore: the solid's motion, the potential,
and one monotone map per fluid region, coupled through the level
geometry and the wall pressure. No fluid displacement is defined,
because none is needed: the state of the fluid is which label sits at
which level.

\subsection{The dictionary to the linear closure}
\label{sec:NL-dictionary}

\emph{Derived.} Linearise about an equilibrium with background
profiles $\hat S_0(c)$, $\hat p_0(c)$ and potential $\Phi_0$.
\begin{itemize}
\item The Eulerian pressure perturbation at a fixed point is
  $p^1 = \delta\hat p(c) + \hat p_0'(c)\,\phi = \delta\hat p - \rho\phi$,
  so $f = p^1 + \rho\phi = \delta\hat p(c)$: the free function of
  \eqref{eq:G-f} is the perturbation of the one-dimensional pressure
  profile at fixed level, which is why it is a function of $\Phi_0$.
\item The label surface of label $S$ sits at level $\hat\Phi_0(S)$,
  the inverse map; its perturbation is the Lagrangian change of the
  potential following the parcels, $\delta\hat\Phi = \Delta\Phi = -s$:
  the lift of \eqref{eq:G-slaving} is minus the perturbation of the
  level map. The slaved normal displacement $u_\nu = -(\phi + s)/g$ is
  then simply the statement that the parcels sit where their label
  sits.
\item The mass below the level value $c$ of the perturbed potential is
  $M_0(c) + \delta M(c)$ with
  $\delta M = \int_{R_c}\rho^1\dV - \rho\bar\Phi$, the second term
  because the region $\{\Phi_0 + \phi < c\}$ has moved by $-\phi/g$.
  With $U = -(\bar\Phi + sA)$ this is
  $\delta M = \int_{R_c}\rho^1 + \rho U + \rho A s$. The linearisation
  of \eqref{eq:NL-mass} is $\delta M = m'(\hat S_0)\,\delta\hat S =
  M_0'\,s = \rho A s$, since the label at level $c$ shifts by
  $\hat S_0'(c)\,s$. The two agree exactly when
  $\int_{R_c}\rho^1 + \rho U = 0$, which is \eqref{eq:G-band} on a
  level wall. The per-level condition of the linear closure is the
  linearisation of the rearrangement constraint, and $V$ is the
  linearised cumulative volume.
\item The fill condition $M(c_{\mathrm{top}}) = m(S_{\max})$
  linearises to \eqref{eq:wall}, and the pressure constant it fixes
  is $f$'s constant.
\end{itemize}
Every object of \S\ref{sec:closure} is thus the first variation of a
named object of the exact theory, and the surface-mass sheet of
\S\ref{sec:dahlen} is the first variation of ``the fluid fills the
container'' (\S\ref{sec:NL-finite}).

\subsection{The second variation, and what the closure computes}
\label{sec:NL-second}

\emph{Argued.} The second variation of \eqref{eq:NL-E} at a secular
state, over the solid's variations and the fluid's class variations
$\boldsymbol\xi$, has the classical form in the fluid (with
$\phi$ the potential perturbation):
\begin{equation}
  E_2^{\mathrm{fluid}} = \frac12\int\kappa\Bigl(\Div\boldsymbol\xi
  - \frac{\rho g\,\xi_\nu}{\kappa}\Bigr)^2\dV
  + \frac12\int\rho N^2\xi_\nu^2\dV
  + \int\rho\,\boldsymbol\xi\cdot\nabla\phi\dV
  + \frac{1}{8\pi G}\int|\nabla\phi|^2\dV ,
  \label{eq:NL-E2}
\end{equation}
the acoustic square, the buoyancy term, the gravitational coupling and
the potential's energy; the along-level part of $\boldsymbol\xi$
enters only through the first and third terms, and the
divergence-free along-level fields (the strict class) not at all. Two
remarks follow, and together they explain the linear closure.

First, the rearrangement class \emph{contains} the overturning
motions: exchanging two parcels of different label is a rearrangement.
So the critical point is a minimiser of the energy on the class exactly
when $N^2 \ge 0$ throughout, and a saddle wherever $N^2 < 0$. The
``secular state'' of an unstably stratified fluid is the continuation
of an unstable equilibrium: a well-defined critical point that the
dynamics does not reach. (The Rayleigh--Taylor growth that the
gauge-penalty treatment shows at $N^2 < 0$ is the dynamics finding the
descent direction; the closure, which imposes equilibrium level by
level, does not look for one.)

Second, the linear closure is the \emph{Schur complement} of
\eqref{eq:NL-E2} over the pointwise normal sector. At fixed $(\phi, s,
\bu_s)$ the stationarity of \eqref{eq:NL-E2} in the pointwise normal
displacement is the slaving $\xi_\nu = -(\phi + s)/g$, exact whenever
$N^2 \ne 0$ of either sign; eliminating $\xi_\nu$ at its stationary
value and the along-level displacement at its (energetically neutral)
compatibility leaves a functional of $(\phi, s, V, \bu_s)$, which we
claim is \eqref{eq:L}. This is the mechanism behind the fact, observed
in \S\ref{sec:closure}, that the sign of $N^2$ never enters: the
eliminated block is positive for $N^2 > 0$ and negative for $N^2 < 0$,
but the reduced expression is the same algebraic object in both cases.
The price is stated above: for $N^2 < 0$ the reduced functional can be
coercive, and the closure unique, while the state it computes is a
saddle of the full energy. At $N^2 = 0$ the eliminated block is
singular, the slaving degenerates, and $\xi_\nu$ joins the kernel ---
class K --- which is the other dichotomy of
\code{doc/static\_fluid\_core.tex}; the level-mean sector keeps its
compressional stiffness, which is why the closure stays well posed
there. The identification of the Schur complement with \eqref{eq:L} is
the step not written out; it would make the symmetry and the
coercivity statements of \S\ref{sec:settled} consequences rather than
observations. We note that an earlier formulation of this point ---
that the ordering of parcels is fixed within the class, so that
overturning is not a variation --- is wrong, and is withdrawn here.

\subsection{Finite-amplitude features}
\label{sec:NL-finite}

\emph{Argued.} Three features of the linear closure that look singular
or accidental are regular at finite amplitude.

\emph{The wall.} The fluid fills the deformed container. The deformed
wall is never exactly level, so near the top of a fluid region the
level sets of $\Phi$ cut the wall over a range of levels of the order
of the deformation, and that range is filled by the lightest fluid
spread along the wall, as a transverse wall is in
\S\ref{sec:transverse}. There is no sheet and no surface mass; the
surface mass of F3 is the first derivative of this picture at a level
wall, where the cut range collapses to the wall's level and its
content to a surface density. The linear closure's level-limit
singularity ($1/\sin\theta$ in the along-level displacement) is the
derivative of a regular object taken at the point where the object's
support degenerates.

\emph{Non-monotone labels.} Where $N^2$ changes sign inside a region
($\hat S_0$ has an extremum), the map $\hat S$ folds and
\eqref{eq:NL-mass} must be read with the measure: $\mu$ is the
push-forward of $\rho\dV$ under $\hat S\circ\Phi$, and the mass with
labels in $E$ is a sum over the branches of $\hat S^{-1}(E)$. The
one-dimensional unknown is then better taken as the pressure profile
$\hat p(c)$ with $\rho(c) = \rho(\hat p(c), \hat S(c))$ and the label
assigned by cumulative mass, which is the variable the linear closure
uses ($f$ and $V$ are a pressure and a cumulative volume, not labels).
The linear closure is regular across $N^2 = 0$ for this reason.

\emph{The foliation.} At finite amplitude the level foliation is that
of the deformed body's geopotential, and the questions of
\S\ref{sec:setting} are asked of it; the convexity results apply to
the deformed region, and the assertion of \S\ref{sec:implementation}
is made on the current iterate.

\subsection{The nonlinear computation}
\label{sec:NL-compute}

\emph{A programme.} The unknowns are the solid's finite motion (the
referential class's variables, with $\varphi_e$ now the current
iterate), the potential, and one one-dimensional profile per fluid
region. The natural one-dimensional coordinate is the volume fraction
$v$ of \S\ref{sec:implementation} read as a Lagrangian coordinate:
with the profile written as the level $\hat\Phi(v)$ at which the
fraction $v$ of the region's volume lies below, the rearrangement
constraint is the statement that the volume of $\{\Phi < \hat\Phi(v)\}$
is $v|\Omega_F|$, computed by the same histogram that the linear
closure uses. The coupled problem is then: the solid's finite
elasticity with the follower load $\hat p(\Phi)$ on the wall; Poisson's
equation with the fluid's density $\rho(\hat p, \hat S)$ on the level
foliation; and the profile equations \eqref{eq:NL-profile}--%
\eqref{eq:NL-mass} on the one-dimensional grid. Its Newton step is the
bordered system of \S\ref{sec:weak} assembled on the current iterate,
with the level geometry $A(c)$ and the wall levels re-read from the
current potential at every step. Nothing in the fluid is a field on a
mesh; the fluid's whole state is a few hundred numbers.

Three relations place this programme. The wholly fluid body (no solid)
is the hydrostatic figure of \code{doc/equilibrium\_figures.tex}, and
its existence theory is that of self-gravitating fluid masses as
energy minimisers at fixed mass and, for rotation, angular momentum
(Auchmuty and Beals; Lieb and Yau for the non-rotating case), with the
label structure adding the rearrangement constraint of Burton; the
nonlinear closure is a hydrostatic figure inside a deformable elastic
container, and the two programmes should share their fluid machinery.
With $\bOmega$ dynamical, the exact angular momentum carries the fluid
only through $\int\rho\,\bx\times(\bOmega\times\bx)$, there being no
relative fluid motion in a secular state, which is the nonlinear form
of the statement in \S\ref{sec:rotation} that the solid-restricted
Tisserand gauge is forced. And the nonlinear gauge and Maxwell routes
of the library carry a fluid displacement, finite-elastic with a gauge
or viscoelastically relaxed, whose nonlinear content is a trajectory
through the rearrangement class; this route carries none, and is for
that reason the natural nonlinear statement of the secular state,
where those routes are the natural statements of how it is reached.

\emph{Open.} The identification of the reduced second variation with
\eqref{eq:L}; existence of the secular state for a stratified fluid in
a deformable container (the whole-fluid and solid-only cases exist in
the literature, their coupling does not); the regularity of the level
map at folds; and the foliation question at finite amplitude.

%=====================================================================
\section{Implementation}
\label{sec:implementation}

\emph{Status: the one-dimensional implementation is done
(\S\ref{sec:verification}); the three-dimensional one is specified
here.}

\paragraph{The coordinate.} The border lives on a one-dimensional grid
in the volume-fraction coordinate
\begin{equation}
  v = \frac{|\{\Phi_0 < c\}\cap F|}{|F|},
  \qquad
  \mathrm dv = \frac{A\dc}{|F|},
  \qquad
  \int\frac{\kappa V'^2}{A}\dc = \frac{1}{|F|}\int\kappa\Bigl(\frac{\mathrm dV}{\mathrm dv}\Bigr)^2\mathrm dv ,
\end{equation}
in which the stiffness has no singularity at a centre ($A \to 0$),
$V(0) = 0$ is a plain essential condition, and $v(\bx)$ is the
cumulative histogram of $\Phi_0$ over the fluid quadrature points
($v = (r/c)^3$ on the sphere, $(r/c)^2$ on the disc). The grid is
coarse against the mesh's resolution of $\Phi_0$, so that every
one-dimensional element contains many quadrature points. The
foliation assumption of \S\ref{sec:setting} is asserted at setup:
$|\nabla\Phi_0|$ bounded away from zero at every fluid quadrature
point outside a small ball about a centre, and one connected component
per level, read off the fluid element graph thresholded at a few
values of $c$; both are cheap and both refuse rather than mislead.

\paragraph{Every coupling is an ordinary integral.} By
\eqref{eq:coarea}, $\int\bar X(c)\,q(c)\dc = \int_F q(\Phi_0(\bx))\,\chi(\bx)\dV$
for any one-dimensional $q$. The columns of \eqref{eq:rows} are
therefore volume integrals over the fluid elements with the border's
basis functions, or their derivatives, composed with $\Phi_0$ as
coefficients --- assembled where F1 is assembled now --- and wall
integrals with the same coefficients. The one-dimensional block is
sparse. $A(c)$, where it is needed pointwise, is the same histogram
differentiated.

\paragraph{Unknowns, as built (\code{SetLevelBorder}).} The mixed
class carries, per fluid region, the nodal unknowns $(U_k, s_k, \nu_k)$
on the grid and one multiplier per level wall, in the form of the
Lagrangian
\begin{equation}
  L_{\mathrm{border}} = \frac12\int\frac{\kappa U'^2}{A}\dc
  + \int\rho\,s\,U'\dc - \frac12\int\rho_\Phi A\,s^2\dc
  + \int\nu\,(sA + U + \bar\Phi)\dc
  + \sum_{\text{walls } j}\lambda_j\Bigl(U(c_j) - \oint_{\Sigma_j}\bu_s\cdot\nuh\dS\Bigr),
  \label{eq:L-built}
\end{equation}
which is \eqref{eq:L} with the lift $s$ kept as an unknown and
$\nu = -f'$ as the multiplier of its definition. Its rows reproduce
\eqref{eq:V-ode} (eliminate $\nu$ between the $s$- and $U$-rows), the
Poisson row receives $+\int\nu\,\chi\dV$, the extra density $-f'$, and
each wall row receives the wall pressure $\lambda_j = \pm f(c_j)$. The
advantage over \eqref{eq:L} is locality: the potential enters only
through the $\nu$ constraint, as the columns $\int N_k(\Phi_0)\chi\dV$,
and the solid only through the wall columns
$-\oint(\mathbf m\cdot\nabla\Phi_0/|\nabla\Phi_0|)(\mathbf m\cdot\bv)\dS$,
so every other border unknown has a zero column and the shared
elimination engine (\code{bordered.hpp}) skips it: the cost is one base
solve per $\nu_k$ and per wall, cached across solves at a frozen
operator. The one-dimensional block uses the element-mean coefficients
of the grid, $\kappa_e$, $\rho_e$ and the element volumes, so that
$\int\kappa U'^2/A\dc$ becomes $\sum_e\kappa_e(\Delta U_e)^2/\mathrm{Vol}_e$
with no $A$ anywhere (and a factor $4/3$ on a centre's first element,
where $A \propto \sqrt{c - c_{\min}}$); the $\nu$ columns are assembled
on the block's own quadrature points, so that the discrete mass balance
of the border is an identity. On a shell the inner wall carries its own
constraint and multiplier. The monolithic MINRES route is not needed at
this width (ten to thirty potential-coupled columns) and is refused
with the border for now, as is the centrifugal coupling of
\S\ref{sec:rotation}, whose columns are not yet wired. For a rotating
figure the geopotential is passed as the background potential.

\paragraph{The node count.} The grid resolves the profile, so too few
nodes under-resolve it; but each node's hat function is integrated over
the mesh's quadrature points, so a grid finer than the mesh's resolution
of $\Phi_0$ makes the columns noisy. On the \code{fluid\_core} cases
below the answer drifts with the count above about twenty nodes at
$h = 0.3$ and $0.2$, and twelve to sixteen nodes are right for those
meshes; the default is sixteen. A rule tied to the mesh (one node per
few element-widths of $\Phi_0$) is the natural refinement.

\paragraph{Cost.} One symmetric solve plus the cached column solves.
No relaxation, no time-stepping, no fluid displacement, no gauge
penalty. The column solves are the whole price: one base solve per
level node (the $\nu$ column) and per level wall, with the operator,
the preconditioner and the iteration counts of the base solve
unchanged, and the condensed border system dense but of order the node
count. Measured on the \code{fluid\_core} case of \S\ref{sec:verify}
at $h = 0.2$, order two, eight ranks: the degree-0 solve takes
$0.54\,$s without the border and $7.5\,$s with sixteen nodes
($18$ columns) or $13.6\,$s with thirty-two, and the degree-1 and
degree-2 solves that follow at the same operator take $0.5\,$s either
way, the columns being cached. In a Love sweep or an elastic time
series the border is therefore paid once; a stepper that changes the
operator pays it once per operator version (the engine keeps two
versions, which covers a stepper alternating between relaxed and
unrelaxed operators). Should the one-off cost matter, the levers are a
multiple-right-hand-side solve for the columns, which share the
operator, and the monolithic route.

\paragraph{The referential variant (a plan).} In the referential class
(\S\ref{sec:solid}) the pieces are, in order of construction: (i)
eliminated fluid regions in that class at all, with the solid's
SubMesh excluding them and the fluid's density on the parent, as the
mixed class has; (ii) the potential glued by the prescribed radial
extension of the slipping classes, $\phi^1 = \zeta^1 - (E\bu)\cdot\nabla\Phi_0$
in the fluid, so that F1, the fluid half of F3 and the $\nu$ columns
become forms in $(\zeta, \bu)$ through $E$ (the existing SubMesh
pairing matrices carry $E$ to true dofs); (iii) F2 and the follower
kernel on the wall, the latter the one-sided solid kernel of the
slipping classes' interface pressure form at $\pi = p^0$; (iv) the
border itself, unchanged in content, assembled on the reference fluid
mesh with the mapping's Jacobians and the wall integrals with the
surface Jacobian; (v) the consistency check $\Se\bn = -p^0\bn$ on the
wall. The first test is the hydrostatic natural state, where the result
must equal the mixed class's to solver tolerance; the second a relabelled
sphere through the mapping, where it must equal the unmapped one; the
third a pre-stressed aspherical solid. None of this is started. The
code map: the welded referential class carries the prescribed vacuum
extension and the mapped Poisson operator but no fluid region of any
kind; the radial fluid extension (\code{NewRadialFluidExtension}, one
fluid core inside a solid shell, with \code{SetFluidExtension}) and the
broken-$\zeta$ door live on the slipping class; the one-sided interface
pressure kernel is \code{SlipInterfacePressureIntegrator}. Step (i)
therefore adds \code{FluidRegion}s to the welded referential class on
the mixed class's pattern, step (ii) lifts the extension from the
slipping class, and the rest is assembly.

%=====================================================================
\section{Verification}
\label{sec:verify}
\label{sec:verification}

\paragraph{The one-dimensional prototype (measured).} The radial tool
of the disc family (\code{benchmarks/disc/disc\_radial.py}, method
\code{dahlen} at $l = 0$) assembles \eqref{eq:rows} from the
level-foliation equations --- the border as $V = r\,u_f$ on the core
nodes, the multiplier as a separate unknown, F1--F3 kept and their
level means cancelled by the border's terms --- and compares with the
welded compressible solve, which at $l = 0$ is exact physics. On a
uniform core ($N^2 < 0$), the neutral core ($N^2 = 0$) and a stably
stratified twin ($N^2 > 0$):
\begin{center}
\small
\begin{tabular}{@{}ll@{}}
\toprule
quantity, closure against welded & worst of the three cores \\
\midrule
surface and CMB displacement & $10^{-11}$ \\
fluid degree-0 profile & $1.4\times10^{-11}$ \\
surface potential, through $\zeta = g\,u_r + \phi$ & $5\times10^{-12}$ \\
multiplier identity $\lambda = -\kappa\Div\bu + \rho_F\phi$ at the wall & $2\times10^{-12}$ \\
resolution drift of the closure, $p = 8 \to 11$ & $2\times10^{-12}$ \\
\bottomrule
\end{tabular}
\end{center}
The reduction alone (F1--F3, no border) gives $u_r(a) = +3.29$ against
$-0.0991$ on the uniform core, and the same number on the neutral core
and its stable twin: the arbitrary family member, $\kappa$-blind as
\eqref{eq:G-f} says it must be. These are the validation entries
``border l=0'' of the tool.

\paragraph{The three-dimensional gate (measured).} The mixed class
on the \code{fluid\_core} Love-number cases (uniform core, so
$\rho_\Phi = 0$ and the closure's mass balance is exact), order 2,
\code{love\_benchmark -level-border}, against the \code{pyslfp}
reference $h'_0 = -0.13210$, $k'_0 = 0$:
\begin{center}
\small
\begin{tabular}{@{}lcccccc@{}}
\toprule
$h'_0$ by level nodes & 8 & 12 & 16 & 24 & 32 & reduction alone \\
\midrule
$h = 0.3$ (3733 elements) & $-0.1306$ & $-0.1309$ & $-0.1310$ & $-0.1326$ & $-0.1340$ & $-1.457$ \\
$h = 0.2$ (8183 elements) & $-0.1311$ & $-0.1314$ & $-0.1315$ & $-0.1328$ & $-0.1337$ & $-1.424$ \\
\bottomrule
\end{tabular}
\end{center}
At sixteen nodes the error is $0.8\,\%$ at $h = 0.3$ and $0.45\,\%$ at
$h = 0.2$, the level at which degrees 1 and 2 themselves sit on these
meshes ($0.2$ to $0.7\,\%$), and it falls with $h$; $k'_0$ is
$-9\times10^{-7}$ and $-7\times10^{-7}$ (against $-5.3$ and $-5.2$ for
the reduction alone), mass conservation to solver tolerance; and the
degree-1 and degree-2 numbers with the border equal those without it
to $4\times10^{-5}$ at $h = 0.2$, the border being inert there. On the
canned three-layer test mesh (a fluid shell, two walls, stratified
density) the closure agrees with the gauged Adams--Williamson twin
under a uniform load to $0.4\,\%$ in the potential and $2.7\,\%$ in
displacement at order 2, from $4.7\,\%$ and $13\,\%$ at order 1
(\code{TestMixedProblemGauged}, the \code{LevelBorder} cases), where the
reduction alone differs by more than $10^{-2}$ by construction.

\paragraph{Gates remaining, in order.}
\begin{enumerate}
\item The slaved fluid trace against the solid's at the CMB on the
  sphere at $\ell = 0$ (pointwise, not only in the level mean), and
  the $\ell \ge 1$ inertness at $10^{-10}$ rather than $10^{-5}$
  (the discrete level sets are spheres only to mesh accuracy).
\item The rotational measurement of \S\ref{sec:rotation}: the
  $\omega_3$ response under a load with a degree-0 part on a
  fluid-core model, with and without the border, against the spectral
  reference.
\item A neutral aspherical equilibrium figure, on which the closure
  must equal the welded and gauged answers.
\item A non-neutral figure: $h$- and $p$-convergence, and a
  first-order-in-flattening check against the spherical response.
\item The rotating figure with the border against the input-inertia
  rotational door on the sphere as control.
\end{enumerate}

%=====================================================================
\section{Standing of the method}
\label{sec:standing}

For static problems with hydrostatic fluid regions the closure, once
the gates above are passed, replaces the Maxwell relaxation of
\code{doc/static\_fluid\_core.tex}: one symmetric solve for the secular
state on any geometry, any stratification and any number of regions,
with the fluid's bulk modulus consulted exactly where the physics
consults it (the level-mean sector) and nowhere else. Two limits
stand. The closure does not touch $\ell \ge 1$ on the sphere, where it
is Dahlen's method; whether Dahlen's degree-1 answer or the welded one
is the secular state on a stably stratified core is decided by the
adjudication in that note, not here. And it is static: the relaxation
remains the only route to the transient, and the dynamic check of the
closure's answer. The gauge-penalty and slipping formulations are
instruments; their standing is unchanged.

\paragraph{Open items.} The coercivity of \eqref{eq:L} off the sphere.
The term-by-term symmetry of the follower load with a deviatoric
$\Se$ on a transverse wall, and the pull-back of the border under
$\varphi_e \ne \mathrm{id}$. The gauge-shift certificate with the
Eulerian fluid on an aspherical level wall. The sufficiency of the
per-level condition \eqref{eq:G-band} on a transverse wall (one
condition per component is the right count; the existence of the
along-level field on an open level surface with its Neumann data is
standard, but is not written out).

\section*{References cited for the foliation}

\begin{itemize}\small
\item L.~G. Makar-Limanov, Solution of Dirichlet's problem for the
  equation $\Delta u = -1$ in a convex region, Math. Notes 9 (1971).
\item A.~U. Kennington, Power concavity and boundary value problems,
  Indiana Univ. Math. J. 34 (1985) 687--704; Theorem 4.2 as restated
  in K. Ishige and P. Salani, Parabolic power concavity and parabolic
  boundary value problems, Math. Ann. 358 (2014), Theorem 1.1.
\item L. Caffarelli and J. Spruck, Convexity properties of solutions
  to some classical variational problems, Comm. PDE 7 (1982)
  1337--1379 (the convex-ring theorem; also R. Gabriel 1957 and
  J. Lewis 1977 for the harmonic and $p$-harmonic cases).
\item F. Hamel, N. Nadirashvili and Y. Sire, Convexity of level sets
  for elliptic problems in convex domains or convex rings: two
  counterexamples, Amer. J. Math. 138 (2016) 499--527.
\item B. Gidas, W.-M. Ni and L. Nirenberg, Symmetry and related
  properties via the maximum principle, Comm. Math. Phys. 68 (1979).
\item B. Kawohl, Rearrangements and convexity of level sets in PDE,
  Lecture Notes in Mathematics 1150, Springer (1985).
\end{itemize}

\section*{References cited for the nonlinear theory}

\begin{itemize}\small
\item G.~R. Burton, Rearrangements of functions, maximization of
  convex functionals, and vortex rings, Math. Ann. 276 (1987)
  225--253, and the steady-flow papers that followed.
\item J.~F.~G. Auchmuty and R. Beals, Variational solutions of some
  nonlinear free boundary problems, Arch. Rational Mech. Anal. 43
  (1971) 255--271 (rotating self-gravitating fluids as energy
  minimisers).
\item E.~H. Lieb and H.-T. Yau, The Chandrasekhar theory of stellar
  collapse as the limit of quantum mechanics, Comm. Math. Phys. 112
  (1987) 147--174 (the non-rotating minimisation, in its classical
  part).
\item T.~B. Benjamin, The alliance of practical and analytical
  insights into the nonlinear problems of fluid mechanics, Lecture
  Notes in Mathematics 503 (1976) 8--29 (steady flows as extremisers
  over rearrangements).
\end{itemize}

\end{document}
